\documentclass[10pt,a4paper]{amsart}
\usepackage[margin=19mm]{geometry}
\usepackage{amsmath,amssymb,mathtools}
\usepackage[scr=rsfs,cal=euler]{mathalfa}
\usepackage{xfrac}
\usepackage[unicode,hidelinks,bookmarks=false]{hyperref}
\newcommand{\ie}{i.e.\ }
\newcommand{\cf}{cf.\ }
\newcommand{\resp}{respectively\ }
\newcommand{\Subsection}{\subsection}
\providecommand{\nobreakdash}{\nobreak\hbox{-}}
\setlength{\emergencystretch}{2em}
\multlinegap0pt
\usepackage[shortlabels]{enumitem}
\usepackage{amssymb}
\usepackage{mathtools}
\usepackage{float}
\usepackage{tikz}
\newtheorem{lemma}{Lemma}
\title{The Combinatorics of Multi-Lane Merging}
\author{Aurora Hiveley, Doron Zeilberger}
\date{Source: arXiv:2607.26334v1 (CC BY 4.0)}
\begin{document}
\maketitle

\noindent\emph{Source and licence.} The Combinatorics of Multi-Lane Merging. Version arXiv:2607.26334v1 (CC BY 4.0), distributed by arXiv under the Creative Commons Attribution 4.0 International licence, http://creativecommons.org/licenses/by/4.0/, as recorded in the arXiv licence metadata for this exact version and shown on the abstract page https://arxiv.org/abs/2607.26334v1. Full licence notice: \texttt{licenses/arxiv-2607.26334-CC-BY-4.0.txt}.\par\smallskip

\noindent\textit{Editorial scope.} Counted unit: the article's lemma lem:bounce (source0.tex lines 250--258). The article's complete original proof of it is delivered with it (source0.tex lines 260--266, one \texttt{proof} environment of 7 lines). No mathematics is written, repaired or re-derived: the statement and the proof are the article's own lines, in the article's own notation, and no retained line is substituted or re-flowed.  No further source-local context is needed: every cross-reference of every retained line resolves inside the two delivered spans.  Nothing was omitted from any retained span, so the package carries no omission record.  No retained line contains a citation, so the wrapper carries no bibliography and no bibliography entry.  No retained line contains a figure environment, an image-inclusion command or drawing code, so no image is delivered and the package declares no asset dependency.  Editorial formatting only, in addition: an AMS \texttt{amsart} preamble that restates the article's own declarations and macros used by the retained lines, namely the environments \texttt{lemma} on separate counters, together with the standard packages those lines need.  The retained environments print in this document as Lemma 1 in order of appearance, whereas the article prints the same items with the numbering of its own full document.  The complete native source of the article as distributed in the arXiv e-print for this exact version is bundled unchanged as \texttt{sources/206260/source0.tex}.
\begin{lemma} \label{lem:bounce}
    If $T$ is a SYT and $T'$ is the SYT obtained by deleting the largest element, $n$, from cell $c$ of $T$, then
\[ \frac{\#A(T)}{\#A(T')} = \beta(c)\]
where
    $$\beta(c) := \begin{cases}
    \# \left\{ c' = [i', j-1] \mid i' > i, T[c'] < T[c]  \right\} & j > 1 \\
    k - i & j = 1
    \end{cases}$$
\end{lemma}

\begin{proof}
    Note that $\#A(T)/\#A(T')$ is equal to the number of possible preferences that could be in position $n$ of the arrival sequence that would produce this SYT. Let $n$ be in cell $[i,j]$ of tableau $T$, meaning that Car $n$ joined lane $i$ in the arrival sequence for $1 \leq i \leq k$. It suffices to count the number of possible preferences could Car $n$ have had which cause it to join lane $i$.
    
    Certainly if Car $n$'s preference is any digit smaller than $i$, then lane $i$ would never be considered by the car, so we can narrow our search space to $\{i, i+1, \dots, k\}$. However, a preference of $i+1$ will only cause Car $n$ to join lane $i$ if the lanes $i$ and $i+1$ have tied lane lengths. In the tableau $T$, this means that there are already digits occupying rows $i$ and $i+1$ in the column to the left of the element $n$. Extrapolating to lower rows, we see that a preference of $i + r$ will only cause Car $n$ to join lane $i$ if \textit{all} lanes $\{i, i+1, \dots, i+r\}$ have tied lengths, so all rows $\{i, i+1, \dots, i+r\}$ of $T$ have entries in the column to the left of the box containing $n$. Then the number of possible preferences that Car $n$ could have which would cause it to join lane $i$ is equal to $\beta(c)$ where $c = [i,j]$ is the location of Car $n$. In this sense, $\beta(c)$ counts the number of dimensions that the merging path would ``bounce" in, or the number of possible digits that Car $n$ could have that would still cause it to enter lane $i$ and thus fill location $[i,j]$ in $T$. Once we have determined the number of possible digits in position $n$ of an arrival sequence which resolves to $T$, we can ``peel off" $n = T[c]$ from $T$, and consider all possible preferences that Car $n-1$ can have in the same fashion. 
    
    Note, however, that any entry in the first column of $T$ has no column to the left to compare to. But the first car to join any lane can prefer any lane to the left, and it will still join lane $i$. For example, as noted before, Car 1 can prefer any of the $k$ lanes, and it will always join the first lane. If Car 2 joins the second lane, then it could have preferred any of the $k-1$ open lanes to the left, and it would still join lane 2. In general, a car in the first column and the $i$-th row could have $k-i$ preferences corresponding to the empty lanes to the left. Hence for $c = [i,j]$ when $j=1$, $\beta(c)$ is defined differently, as if the zero-th column is full.
\end{proof}

\end{document}
